\section{Proof Sketch}\label{sec:proof-sketch}

The quotient map preserves every finite factor labeling pattern and
every finite Littlestone tree, while the full Cartesian restriction map
preserves independent factor choices.  The structural results in
Proposition~\ref{prop:step-001-structural-certificate} therefore give
the exact quotient dimensions, the product VC and Littlestone
identities, measurable decoded outputs, and the risk identity
\eqref{eq:prelim-exact-risk}.

For the upper bound, each countable quotient supports a fully totalized
measurable VC-one factor kernel.  A common random permutation turns a
one-record change into at most one factor-summary replacement, and
Proposition~\ref{prop:step-003-joint-composition} shows that routing
affects at most two factors.  Their two half-budget privacy costs compose
to \((\varepsilon,\delta)\), including on nonrealizable inputs.
The utility proof does not ask every factor to receive its quota.
Instead, Proposition~\ref{prop:step-004-weighted-shortage} controls the
total distribution mass of short factors, and
Proposition~\ref{prop:step-005-pac-closure} combines this with the
unpadded factor guarantee to obtain the exact global PAC event.
Proposition~\ref{prop:step-006-public-quota-bridge} then sums the
heterogeneous quotas and absorbs every ceiling into the public rate.

For the lower bound, the additive VC obstruction first pays for the
total mass of factors with small iterated-log complexity.  Each
remaining factor has a subcritical exact budget and satisfies the
one-factor delta cap.  The compact finite game in
Proposition~\ref{prop:step-010-hard-prior} fixes a hard task prior
before seeing the global learner.  A one-use simulator embeds each
factor input row into at most one global row and is therefore private on
all factor inputs.  The only discrepancy between this simulator and the
ideal product experiment is sample-count overflow, whose marginal
probability is bounded by
\(\eta_0=e^7(2/9)^9<3/2048\).

All active factor risks are marginals of one common product-prior
experiment.  Lemma~\ref{lem:step-013-exact-risk-sum} gives an exact
pointwise weighted risk identity, and
Proposition~\ref{prop:step-013-tensorization} charges each overflow
residual once by its factor weight.  The resulting expected global risk
is separated from the PAC expectation ceiling by the exact rational
calculation in Lemma~\ref{lem:step-014-rational-gap}.  Finite
conditioning then selects one deterministic full-product target and
legal mixture with strict PAC failure, yielding
Proposition~\ref{prop:step-014-candidate-closure}.

Finally, Proposition~\ref{prop:step-015-global-upper} establishes the
arbitrary-\(\delta\) upper theorem and finiteness of the least successful
sample size.  Proposition~\ref{prop:step-015-candidate-lower} retains
both lower delta conditions at one fixed candidate.
Proposition~\ref{prop:step-015-conditional-sandwich} combines the two
only after those conditions are checked at the attained sample
complexity.  Propositions~\ref{prop:step-015-one-factor-upper} and
\ref{prop:step-015-one-factor-lower} separately verify the unpadded
upper rule and zero-overflow unrestricted lower baseline when \(k=1\).
